% p274_1; source PDF page 274; manually verified editable table.
\begin{tikzpicture}[x=1cm,y=-1cm,inner sep=0pt,outer sep=0pt]
\path[use as bounding box] (0,0) rectangle (16.000,7.250);
\draw[line width=.4pt] (0.000,0.000) rectangle (3.900,0.850);
\node[align=center,text width=3.640cm,font={\linespread{1}\fontsize{10}{13.50}\selectfont\bfseries}] at (1.950,0.425) {类型};
\draw[line width=.4pt] (3.900,0.000) rectangle (9.200,0.850);
\node[align=center,text width=5.040cm,font={\linespread{1}\fontsize{10}{13.50}\selectfont\bfseries}] at (6.550,0.425) {举例};
\draw[line width=.4pt] (9.200,0.000) rectangle (16.000,0.850);
\node[align=center,text width=6.540cm,font={\linespread{1}\fontsize{10}{13.50}\selectfont\bfseries}] at (12.600,0.425) {结论};
\draw[line width=.4pt] (0.000,0.850) rectangle (3.900,4.050);
\node[align=left,text width=3.640cm,font={\linespread{1}\fontsize{10}{13.50}\selectfont}] at (1.950,2.450) {空间上是现场，但时间上已不是当场};
\draw[line width=.4pt] (3.900,0.850) rectangle (9.200,4.050);
\node[align=left,text width=5.040cm,font={\linespread{1}\fontsize{10}{13.50}\selectfont}] at (6.550,2.450) {甲晚上在乙家盗窃到财物后，在乙家门外草垛睡觉。早上，乙回家发现甲，便抓甲。甲将乙打倒在地};
\draw[line width=.4pt] (9.200,0.850) rectangle (16.000,4.050);
\node[align=left,text width=6.540cm,font={\linespread{1}\fontsize{10}{13.50}\selectfont}] at (12.600,2.450) {空间上，盗窃行为与暴力行为是同一现场，但二者在时间上已经不是当场，不转化为抢劫};
\draw[line width=.4pt] (0.000,4.050) rectangle (3.900,7.250);
\node[align=left,text width=3.640cm,font={\linespread{1}\fontsize{10}{13.50}\selectfont}] at (1.950,5.650) {空间上不是同一现场，但时间上是当场。主要是指持续追捕的情形};
\draw[line width=.4pt] (3.900,4.050) rectangle (9.200,7.250);
\node[align=left,text width=5.040cm,font={\linespread{1}\fontsize{10}{13.50}\selectfont}] at (6.550,5.650) {甲抢夺了乙的手机，迅速逃跑，乙立即追，持续追了1公里后，甲为抗拒抓捕，将乙打倒在地};
\draw[line width=.4pt] (9.200,4.050) rectangle (16.000,7.250);
\node[align=left,text width=6.540cm,font={\linespread{1}\fontsize{10}{13.50}\selectfont}] at (12.600,5.650) {空间上，盗窃行为与暴力行为不是同一现场，但持续追捕，表明二者在时间上具有密切联系，属于当场，能够转化为抢劫（2025年试题）};
\end{tikzpicture}
